\documentclass[UTF8, fontset=macnew, a4paper, 12pt]{ctexart}

\usepackage{geometry}
\geometry{top=2.2cm, bottom=2.2cm, left=2.2cm, right=2.2cm}
\usepackage{amsmath, amssymb}
\usepackage{booktabs}
\usepackage{multirow}
\usepackage{tabularx}
\usepackage{graphicx}
\usepackage{caption}
\usepackage{xcolor}
\usepackage{float}
\usepackage{enumitem}
\usepackage{fancyhdr}
\usepackage{setspace}
\usepackage{pifont}
\usepackage{natbib}
% 拉丁作者名之後採半形括號（中文作者名在書目的中文部分仍用全形）；
% 第五個引數留空表示作者與年份之間不加逗號。
\bibpunct{(}{)}{; }{a}{}{,}
\usepackage[hidelinks]{hyperref}

\graphicspath{{../../output/figures/}}
\captionsetup{font=small, labelfont=bf, labelsep=quad}
\setlist{itemsep=2pt, parsep=0pt, topsep=4pt}
\singlespacing

% ctexart 預設的圖表標號前綴為簡體「图」，改回正體。
\renewcommand{\figurename}{圖}
\renewcommand{\tablename}{表}

\ctexset{
  section/format = \raggedright\zihao{4}\bfseries,
  section/name   = {},
  section/number = \bignum{\value{section}},
  section/aftername = 、,
  subsection/format = \raggedright\zihao{5}\bfseries,
  subsection/number = \chinese{subsection},
  subsection/aftername = 、,
  subsubsection/format = \raggedright\normalsize\bfseries,
}
\newcommand{\bignum}[1]{\ifcase#1\or 壹\or 貳\or 參\or 肆\or 伍\or
  陸\or 柒\or 捌\or 玖\or 拾\fi}
\renewcommand{\thesection}{\bignum{\value{section}}}
% 交叉引用子節時補上節次，否則 \ref 只印出「三」而無從判斷屬於哪一節。
\renewcommand{\thesubsection}{\chinese{subsection}}
\makeatletter
\renewcommand{\p@subsection}{\thesection 節之}
\makeatother
\setcounter{secnumdepth}{2}

\pagestyle{fancy}
\fancyhf{}
\fancyhead[L]{}
\fancyhead[R]{\small 進度報告（1012 理論初探）}
\fancyfoot[C]{\small \thepage}
\renewcommand{\headrulewidth}{0.4pt}

% U+2460~2463 不在 xeCJK 認定的 CJK 範圍內，會落到拉丁字型而缺字，
% 因此改用 pifont 的圈號。
\newcommand{\cn}[1]{\ding{\number\numexpr171+#1\relax}}
\newcommand{\edag}{e_0^{\dagger}}
\newtheorem{prop}{Proposition}
\newtheorem{cor}{Corollary}
\newtheorem{lem}{Lemma}
\newenvironment{proof}{\par\noindent\textit{Proof}.\ }{\hfill$\square$\par\vspace{4pt}}
\newcommand{\Ex}{\mathbb{E}}

% 分部標題
\newcommand{\partline}[2]{%
  \par\vspace{10pt}\noindent\rule{\textwidth}{0.8pt}\par\vspace{2pt}
  \noindent{\zihao{4}\bfseries 第#1部\quad #2}\par
  \vspace{2pt}\noindent\rule{\textwidth}{0.8pt}\par\vspace{6pt}%
  \phantomsection\addcontentsline{toc}{part}{第#1部\quad #2}}

% 目錄：只收到小節，部次以粗體與一行間距與節次區隔。
\setcounter{tocdepth}{2}
\makeatletter
\renewcommand{\l@part}[2]{\par\vspace{8pt}%
  \noindent{\bfseries #1}\par\vspace{2pt}}
\renewcommand{\l@section}[2]{\@dottedtocline{1}{0em}{5.2em}{\bfseries #1}{\bfseries #2}}
\renewcommand{\l@subsection}[2]{\@dottedtocline{2}{5.2em}{4.6em}{#1}{#2}}
\makeatother

\begin{document}
\pagenumbering{roman}

\begin{center}
{\zihao{3}\bfseries 修勻、資訊加權與中心化校正的代數關係}\\[4pt]
{\zihao{5} ——理論初探}\\[6pt]
{\zihao{5} The Algebra Shared by Graduation, Information Weighting,}\\[2pt]
{\zihao{5} and Centring Correction: A First Exploration}\\[12pt]
{\zihao{5} （林子立）\quad 指導教授：余清祥\,教授}\\[4pt]
{\zihao{5} 國立政治大學統計學系}\\[4pt]
{\zihao{5} 中華民國 115 年 10 月 12 日}
\end{center}

\vspace{6pt}
\begin{center}
\begin{minipage}{0.92\textwidth}
\small
\textbf{摘要}\par\vspace{4pt}
本文把三件原先各自成立的工具寫進同一組正規方程式
$\hat r=(W+h\Delta^{\top}\!\Delta)^{-1}W(y-b)$：資訊加權供其中的 $W$，
修勻供 $h\Delta^{\top}\!\Delta$，而對數轉換偏誤的中心化供位移項 $b$。
由此得到五項結果。其一，沿年齡修勻之後，$b(\mu;c_0)$ 不再是該估計量的偏誤，
應扣除的是 $S\,b(\hat\mu)$，其中 $S$ 為修勻算子；原樣扣除的誤設量隨平滑量
增大，實測中 $\alpha_x$ 的最大偏誤由 $0.672$ 升至 $11.072$，而依此改正後
在三十三倍的平滑量全距內只落在 $0.213$ 至 $0.303$。其二，沿時間修勻則
不需要任何輸送：對任一凸的懲罰 $J(\Delta^{z}r)$，$W$ 加權平均恆不變，
而 $z\ge2$ 時 $W$ 加權的最小平方斜率亦恆不變，故 $\hat\alpha_x$ 與漂移項
精確不動，修勻只作用於 $\beta_x$ 與時間參數的形狀。
其三，該不變性來自差分算子的零空間而非二次範數，故把懲罰改為 $L_1$
不影響兩項保證，所改變的是解的類別——由處處彎曲改為分段多項式，
而其貝氏讀法由常態先驗改為 Laplace 先驗，亦即常態的尺度混合。
其四，高階修正有三種可能的對象，其中插入式誤差可分解為
$b'(\mu)\,\mathrm{bias}(\hat\mu)+\tfrac12 b''(\mu)\mathrm{Var}(\hat\mu)$ 兩項，
實測與該分解的逐年齡相關為 $0.817$ 與 $0.899$，但兩項在最關鍵的幼年組
幾乎互相抵銷（人數五萬、$1$--$4$ 歲：$-0.113$ 與 $+0.075$，而實際為 $+0.001$），
故直接施以二階校正在數值上不穩定。
其五，修勻的先決條件可寫成「懲罰的零空間須包含所相信的模型」，
而 $z=2$ 的零空間即對數死亡率在格子上成一次式——
沿年齡時這是 Gompertz 法則、沿時間時這是固定漂移的假設，
兩者皆為本文已有的結構假設。又該條件所需的量全部有閉式，
故可在配適之前逐年齡算出：臺灣女性資料上它把 $5$--$9$ 至 $20$--$24$ 判為惡化、
$25$--$29$ 以上判為改善（均方誤差降至四成至八成），
與對數死亡率曲度的分布相符。
又零格的處理可依列資訊量 $I_x=\sum_t\mu_{x,t}$ 與曲度分三級，
而丙級（須借用外部資訊者）恰為 $0$ 至 $20$--$24$ 的六組，
與均方誤差比所判者完全相同。
\par\vspace{6pt}
\textbf{關鍵詞}：Whittaker 修勻、trend filtering、混合模型表示、
零空間不變性、插入式估計、高階偏誤修正
\end{minipage}
\end{center}

\clearpage
{\zihao{4}\bfseries 目\quad 錄}\par\vspace{8pt}
\makeatletter\@starttoc{toc}\makeatother

\clearpage
\pagenumbering{arabic}

\section{修勻、資訊加權與中心化的代數關係}\label{sec:algebra}

本文先前確立的結果有三。其一，標準 Lee-Carter 模型對
$\log\hat m_{x,t}=\log(\max(D_{x,t},c_0)/E_{x,t})$ 作奇異值分解，
而 $\hat\alpha_x$ 即該列的平均，故其偏誤為逐格偏誤的平均，
且逐格偏誤有閉式 $b(\mu;c_0)=\Ex[\log\max(D,c_0)]-\log\mu$。
其二，該偏誤可事後扣除，此即 \citet{coxsnell1968} 意義下的修正，
與 \citet{firth1993} 的預先修正相對。其三，扣除不處理時間參數的漂移，
須另以期望死亡數加權處理。

\citet{whittaker1923} 的修勻、\citet{lee2003} 的 partial SMR 與
\citet{yue2019} 把兩者接在 Lee-Carter 之前或之後的作法，
在本文的架構中則是另一類改動：它們改動資料而非估計方程。
本節指出三者其實落在同一組正規方程式的三個位置。

\subsection{同一組正規方程式}\label{subsec:one}

Whittaker 修勻解的是

\begin{equation}\label{eq:whit}
\min_r \ (y-r)^{\top} W (y-r) + h\, r^{\top}\Delta^{\top}\!\Delta\, r
\quad\Longrightarrow\quad
(W + h\Delta^{\top}\!\Delta)\,\hat r = W y ,
\end{equation}

其中 $y$ 為觀測的對數死亡率、$W$ 為權重對角矩陣、$\Delta$ 為 $z$ 階差分算子、
$h$ 為平滑量。本文的資訊加權供的正是式~\eqref{eq:whit} 中的 $W$，
修勻供的是 $h\Delta^{\top}\!\Delta$，而中心化供的是把資料端由 $y$ 換成 $y-b$，
於是三者合併為

\begin{equation}\label{eq:unified}
\hat r \;=\; (W + h\Delta^{\top}\!\Delta)^{-1} W\,(y-b),
\qquad S \;:=\; (W + h\Delta^{\top}\!\Delta)^{-1} W .
\end{equation}

式~\eqref{eq:unified} 有兩種讀法，而本文的兩個問題各對應一種。
就計算而言，三個工具落在三個互不相同的位置，故可以疊加；
第\ref{sec:age}與第\ref{sec:time}節檢驗疊加是否真的成立。
就推論而言，$W$ 是概似的精確度而 $h\Delta^{\top}\!\Delta$ 是先驗的精確度，
$\hat r$ 即高斯線性模型的後驗平均，而 $(W+h\Delta^{\top}\!\Delta)^{-1}$ 即後驗變異數。
\citet{kimeldorf1967} 的貝氏修勻給出的正是這個讀法，
\citet{verrall1993} 把它寫成狀態空間的形式，
而 \citet{eilers1996} 的 P-樣條與 \citet{currie2004} 在死亡率上的應用
使 $h$ 成為一個變異成分比，於是可以由 REML 定出
（\citealp{ruppert2003}；\citealp{langbrezger2004}）。

\subsection{兩個方向的修勻並不對稱}\label{subsec:two}

$\Delta$ 可以沿年齡作用，也可以沿時間作用，而兩者的性質截然不同。
沿年齡的修勻是死亡率文獻的主流，\citet{currie2004} 由此處理
$\hat\beta_x$ 在相鄰年齡之間的劇烈起伏，\citet{delwarde2007}
把懲罰直接加在 Poisson 對數概似上，\citet{currie2013} 則給出帶約束與
懲罰的一般框架。沿時間的修勻在死亡率文獻中少見以修勻之名出現，
但標準 Lee-Carter 的外推步驟——$\kappa_t$ 為帶漂移的隨機漫步
\citep{lee1992}——本身即是一個沿時間的平滑先驗，
只是 $h$ 被慣例固定。本文以下兩節分別處理兩個方向，
而兩者對中心化校正的影響並不相同。

\section{年齡方向的修勻與校正量的對應}\label{sec:age}

\subsection{修勻改變估計量，故改變其偏誤}\label{subsec:transport}

沿年齡修勻之後的資料為 $\tilde\ell = S\ell$。由 $\Ex[\ell]=\log\mu+b(\mu)$ 得

\begin{prop}\label{prop:transport}
設 $S$ 為式~\eqref{eq:unified} 的修勻算子，$\tilde\ell=S\ell$。則
$$\Ex[\tilde\ell] \;=\; S\log\mu \;+\; S\,b(\mu),$$
故以 $\log\mu$ 為目標時，$\tilde\ell$ 的偏誤為
$(S-I)\log\mu$ 與 $S\,b(\mu)$ 兩項之和。
\end{prop}

\begin{proof}
$\Ex$ 為線性算子而 $S$ 為常數矩陣，故
$\Ex[S\ell]=S\,\Ex[\ell]=S(\log\mu+b(\mu))$，加減 $\log\mu$ 即得。
\end{proof}

命題~\ref{prop:transport} 的意涵直接：\textbf{修勻改變了估計量， $b(\mu;c_0)$ 因而不再是它的偏誤}。$b$ 是一個特定泛函的偏誤，
即 $\Ex[\log\max(D,c_0)]-\log\mu$；修勻之後的估計量是跨年齡的加權平均，
其偏誤是被 $S$ 重新分配過的 $b$。應扣除的是 $S\,b(\hat\mu)$ 而非 $b(\hat\mu)$，
且原樣扣除的誤設量隨 $h$ 增大。

\subsection{實測}\label{subsec:age-emp}

表~\ref{tab:transport} 以臺灣女性 2001--2024 年的資料驗證。
以全國資料的 Poisson 配適為真值，按該年齡結構縮放至指定人數後
模擬 Poisson 死亡數，重複 100 次，固定種子；
$\lambda=h/\bar w$ 為標準化後的平滑量。
年齡組為 $0$、$1$--$4$、$5$--$9$、\ldots、$100^+$ 共二十二組，
而差分懲罰只施於 $5$--$9$ 歲以上的二十組：前兩組的組距與其餘不同，
二階差分在不等距格子上本即誤設，且嬰幼兒死亡率在年齡上陡降、
其型態本身並不平滑。

\begin{table}[H]
\centering
\caption{校正量輸送前後的對照}
\small
\renewcommand{\arraystretch}{1.1}
\begin{tabular}{llrrrrr}
\toprule
人數 & 估計量 & $\alpha$ 中位 & $\alpha$ 最大 & $\mathrm{SSE}(\beta)$ & $e_0$ 偏誤 & $e_0$ 均方根\\
\midrule
\multirow{7}{*}{五萬}
 & Firth                            & 0.0071 & \textbf{0.0760} & 1.4968 & $-0.091$ & \textbf{0.556}\\
 & 中心化＋加權                      & 0.0235 & 0.0739 & 1.1220 & $+0.203$ & 0.712\\
 & 先修勻、原樣扣 $b$，$\lambda=1$    & 0.0595 & 2.1845 & 1.0118 & $+0.244$ & 0.678\\
 & 先修勻、扣 $Sb$，$\lambda=0.3$     & 0.0381 & 0.3007 & 1.0263 & $+0.185$ & 0.673\\
 & 先修勻、扣 $Sb$，$\lambda=1$       & 0.0390 & 0.2773 & 1.0075 & $+0.213$ & 0.671\\
 & 先修勻、扣 $Sb$，$\lambda=3$       & 0.0551 & 0.2131 & 0.9917 & $+0.197$ & 0.677\\
 & 先修勻、扣 $Sb$，$\lambda=10$      & 0.0406 & 0.3031 & \textbf{0.9508} & $+0.295$ & 0.703\\
\midrule
\multirow{5}{*}{一萬}
 & Firth                            & 0.0129 & \textbf{0.1560} & 20.8896 & $-1.048$ & 1.517\\
 & 中心化＋加權                      & 0.0471 & 0.5064 & 2.6644 & $-1.153$ & 1.802\\
 & 先修勻、原樣扣 $b$，$\lambda=10$   & 0.2256 & 4.1166 & 1.1106 & $-0.241$ & 1.490\\
 & 先修勻、扣 $Sb$，$\lambda=3$       & 0.0963 & 0.3281 & 1.2502 & $-0.195$ & 1.473\\
 & 先修勻、扣 $Sb$，$\lambda=10$      & 0.0878 & 0.3343 & 1.1627 & $\mathbf{-0.001}$ & \textbf{1.480}\\
\bottomrule
\end{tabular}

\vspace{2pt}
\parbox{0.92\textwidth}{\footnotesize 註：種子與本文第伍節不同，
故基準線的數字與該節的表略有差異（人數五萬時中心化加權的 $e_0$ 偏誤
此處為 $+0.203$、該節為 $+0.003$，相差約一點八個蒙地卡羅標準誤），
比較須在同一張表內進行。完整結果見 \texttt{tableG2\_graduation\_transport.csv}。}
\label{tab:transport}
\end{table}

原樣扣除時 $\alpha$ 最大偏誤隨平滑量單調爆開，人數五萬時
$\lambda=0.3,1,3,10$ 分別為 $0.672$、$2.185$、$5.169$、$11.072$，
與命題~\ref{prop:transport} 所預測的「誤設量隨 $h$ 增大」一致。
改扣 $S\,b(\hat\mu)$ 之後爆開完全消失，且在三十三倍的 $\lambda$ 全距內
$\alpha$ 最大偏誤只落在 $0.213$ 至 $0.303$。
此一對平滑量的不敏感性有一項實務意涵：本文已以整節指出零格替代值 $c_0$
的任意性（\citealp{chenroth2024} 在計量經濟學一側指出同一問題），
若在此另引入一個手選的 $h$，論證即不一致；輸送之後 $h$ 幾乎不改動結論，
故該顧慮得以解消，而第\ref{sec:premise}節的 REML 成為加分而非前提。

人數一萬時 $\lambda=10$ 的輸送版在零歲平均餘命上追上 Firth：
偏誤 $-0.001$ 對 $-1.048$、均方根 $1.480$ 對 $1.517$，
而 $\mathrm{SSE}(\beta)$ 為 $1.163$ 對 $20.890$。
平均餘命的離散原是 Firth 唯一全面勝出之處，至此不再成立；
但 Firth 在 $\alpha$ 上仍勝（最大 $0.156$ 對 $0.334$），與本文既有的結論一致。

\subsection{修勻自身的偏誤不宜一併扣除}\label{subsec:nosecond}

命題~\ref{prop:transport} 的另一項 $(S-I)\log\mu$ 為修勻自身的偏誤。
以 $\log\hat\mu$ 代入而一併扣除的版本在實測中發散，
人數一萬、$\lambda=10$ 時零歲平均餘命的偏誤達 $1.8\times10^{6}$ 歲。
成因明確：$\log\hat\mu$ 本身來自修勻後的配適、已經是平滑的，
以它估計修勻的偏誤等於要求迭代自己把平滑反轉回去，每一輪放大一次。
此一現象與非參數推論的既有結論一致——\citet{calonico2018} 指出
為改善涵蓋率而估計並扣除平滑偏誤往往不如直接少平滑一些。
故 $(S-I)\log\mu$ 應留著不動，它是修勻的代價；
$\lambda$ 掃描顯示該代價在本文的規模下很小，
零歲平均餘命的偏誤在 $-0.001$ 至 $+0.295$ 歲之間。

\section{時間方向的修勻與兩個不變性}\label{sec:time}

\subsection{零空間給出的兩項保證}\label{subsec:null}

沿時間方向時，$\hat\alpha_x$ 恰是該年齡各年的 $W$ 加權平均，
而漂移項是時間的斜率。兩者都落在差分算子的零空間附近，
故有以下結果。本節把它寫成對任意凸懲罰皆成立的形式，
蓋因第\ref{sec:l1}節要用到這個一般性。

\begin{prop}\label{prop:invariance}
設 $W$ 為正定對角矩陣、$\Delta^{z}$ 為 $z$ 階差分算子、$J$ 為任一凸函數，
而 $\hat r$ 為
$$\min_r\ \tfrac12 (r-y)^{\top}W(r-y) + J(\Delta^{z} r)$$
的解。記 $\mathbf 1=(1,\ldots,1)^{\top}$、$t=(1,2,\ldots,T)^{\top}$。則
\begin{enumerate}[label=(\roman*)]
\item 對 $z\ge1$，$\mathbf 1^{\top}W\hat r=\mathbf 1^{\top}W y$；
\item 對 $z\ge2$，另有 $t^{\top}W\hat r=t^{\top}W y$。
\end{enumerate}
亦即 $W$ 加權平均恆不變，而 $z\ge2$ 時 $W$ 加權的最小平方斜率亦恆不變。
\end{prop}

\begin{proof}
目標函數為凸，其一階條件為
$W(\hat r-y)+(\Delta^{z})^{\top}g=0$，其中
$g\in\partial J(\Delta^{z}\hat r)$ 為某個次梯度。
左乘 $\mathbf 1^{\top}$ 並利用 $\mathbf 1^{\top}(\Delta^{z})^{\top}=(\Delta^{z}\mathbf 1)^{\top}=0$
（差分算子消滅常數）即得 (i)。
$z\ge2$ 時另有 $\Delta^{z}t=0$（二階以上的差分消滅一次式），
左乘 $t^{\top}$ 即得 (ii)。
加權最小平方斜率為
$\sum_i w_i(t_i-\bar t_w)\hat r_i\big/\sum_i w_i(t_i-\bar t_w)^2$，
其分子等於 $t^{\top}W\hat r-\bar t_w\,\mathbf 1^{\top}W\hat r$，
由 (i) 與 (ii) 知該值不變。
\end{proof}

\begin{cor}\label{cor:alpha}
沿時間以同一組權重 $W$ 施以修勻時，$\hat\alpha_x$ 精確不變；
若另取 $z\ge2$，則 $\kappa_t$ 的加權最小平方斜率亦精確不變。
故沿時間的修勻只作用於 $\beta_x$ 與時間參數的形狀。
\end{cor}

推論~\ref{cor:alpha} 與第\ref{sec:age}節形成明顯的對比。
沿年齡修勻會改動 $\hat\alpha_x$，故必須依命題~\ref{prop:transport} 輸送校正量；
沿時間修勻不動 $\hat\alpha_x$，故與中心化校正完全不互相干擾，
可以直接疊上去而不需任何修正。

\subsection{兩項須注意的事}\label{subsec:time-caveat}

\textbf{其一，漂移估計式必須與懲罰相容，而現行的並不相容。}
本文目前以端點差 $(\kappa_T-\kappa_1)/(T-1)$ 估漂移
（\texttt{R/core/rabbi\_mazzuco\_replication.R} 第 199 行），
而端點差不落在懲罰的零空間，修勻會動它：
在一組 $T=24$ 的核對中由 $0.3179$ 變為 $0.2924$，
而同一組資料的加權最小平方斜率精確不變（差 $5.6\times10^{-17}$）。
改用加權最小平方斜率可使推論~\ref{cor:alpha} 的保證真正生效。

\textbf{其二，推論~\ref{cor:alpha} 的保證對 $\hat\alpha_x$ 是精確的， 對漂移則只是近似。} $\hat\alpha_x$ 就是加權列平均，命題直接適用；
但 $\hat\kappa_t$ 由中心化後加權矩陣的奇異值分解得出，
是資料的非線性函數，故零空間的論證不能逐字套用。
實測顯示漂移的偏誤不因修勻而惡化、在小人口下反而改善，
但這是實證結果而非命題的推論。

\subsection{實測}\label{subsec:time-emp}

\begin{table}[H]
\centering
\caption{沿時間修勻對四項指標的影響}
\small
\renewcommand{\arraystretch}{1.1}
\begin{tabular}{llrrrrr}
\toprule
人數 & $\lambda$ & $\alpha$ 最大 & $\alpha$ 中位 & $\mathrm{SSE}(\beta)$ & 漂移偏誤 & $\kappa$ 粗糙度\\
\midrule
\multirow{4}{*}{一萬}
 & 不修勻 & 0.8046 & 0.1346 & 2.6644 & 0.1850 & 315.97\\
 & 0.1    & 0.8034 & 0.1327 & 2.2584 & 0.1446 & 67.46\\
 & 0.4    & 0.8040 & 0.1299 & \textbf{2.1026} & \textbf{0.1379} & 15.93\\
 & 10     & 0.8053 & 0.1237 & 1.9083 & 0.1288 & 0.31\\
\midrule
\multirow{4}{*}{五萬}
 & 不修勻 & 0.4683 & 0.0663 & 1.1132 & 0.1190 & 88.91\\
 & 0.1    & 0.4700 & 0.0661 & 1.1006 & 0.1201 & 19.81\\
 & 0.4    & 0.4717 & 0.0656 & 1.0999 & 0.1262 & 4.77\\
 & 10     & 0.4742 & 0.0653 & 1.0890 & 0.1282 & 0.08\\
\bottomrule
\end{tabular}

\vspace{2pt}
\parbox{0.92\textwidth}{\footnotesize 註：$\kappa$ 粗糙度為
$\sum_t(\Delta^2\kappa_t)^2$，真值為 $20.57$。$\alpha$ 最大與中位為
各次重複之 $\max_x|\hat\alpha_x-\alpha_x|$ 與
$\mathrm{median}_x|\hat\alpha_x-\alpha_x|$ 的中位數，
故與表~\ref{tab:transport} 的統計量定義不同，兩表的數字不可互相比較。
重複 100 次，固定種子。}
\label{tab:time}
\end{table}

三項讀法。其一，$\alpha$ 最大偏誤在 $\lambda$ 跨一百倍時幾乎不動
（一萬：$0.8046$ 至 $0.8053$；五萬：$0.4683$ 至 $0.4742$），
與推論~\ref{cor:alpha} 一致；殘餘的微小變動來自迭代中 $W=\hat\mu/\bar{\hat\mu}$
隨 $\beta,\kappa$ 更新而略有改變，並非同一組權重。

其二，\textbf{中心化校正把 $\hat\kappa$ 弄得比真值粗糙十五倍}
（$315.97$ 對 $20.57$），而沿時間修勻正好把它修回來；
把 $\lambda$ 校到粗糙度對上真值時，一萬人約在 $\lambda=0.3$、
五萬人約在 $\lambda=0.1$。本文第伍節之四已記錄「校正後的時間路徑比
標準 Lee-Carter 更不平滑」，此處給出該觀察的量化版與其對策。

其三，在校準後的 $\lambda$ 上，人數一萬時 $\mathrm{SSE}(\beta)$ 由
$2.6644$ 降到 $2.1026$、漂移偏誤由 $0.1850$ 降到 $0.1379$，
各約兩成一與兩成七；人數五萬時除 $\kappa$ 的粗糙度以外幾乎沒有變化。
漂移是本文承認中心化碰不到、只能靠加權處理的一項，
故這是第二支獨立的槓桿。

\textbf{但限制要寫清楚：使時間方向安全的那個不變性， 同時使它對 $\alpha$ 完全無用。} $\hat\alpha_x$ 不變是精確的，
故不應期待沿時間的修勻改善 $\alpha$ 的偏誤。
又 $\lambda$ 在時間方向對 $\kappa$ 的形狀影響甚大，
不像年齡方向輸送之後那樣不敏感，故在這個方向
平滑量的估計由加分項回升為必要。

\section{以 $L_1$ 範數取代二次懲罰}\label{sec:l1}

\subsection{兩項不變性不受範數選擇影響}\label{subsec:l1-inv}

把式~\eqref{eq:whit} 的懲罰改為 $h\|\Delta^{z}r\|_1$ 即
\citet{kim2009} 的 trend filtering。命題~\ref{prop:invariance} 已就
任意凸的 $J$ 證明，故 $J(\cdot)=h\|\cdot\|_1$ 是其特例，
兩項不變性原封不動。\textbf{這一點本身即是結論： 不變性來自差分算子的零空間，與二次範數無關。}

表~\ref{tab:l1inv} 以 ADMM 解 $L_1$ 問題後核對。
解法器先以 $h\to\infty$ 核對過：$z=1$ 時退化為加權常數、
$z=2$ 時退化為加權最小平方直線，兩者的最大差為
$7.5\times10^{-10}$ 與$2.9\times10^{-9}$。

\begin{table}[H]
\centering
\caption{$L_1$ 懲罰下兩項不變性的數值核對}
\small
\renewcommand{\arraystretch}{1.1}
\begin{tabular}{llrrr}
\toprule
序列 & $z$ & 加權平均之差 & 加權斜率之差 & $\textstyle\sum_i\mathrm{sign}(\Delta^z\hat r)_i$\\
\midrule
\multirow{2}{*}{雜訊無趨勢}
 & 1 & $5.4\times10^{-16}$ & $2.2\times10^{-2}$ & $-1$\\
 & 2 & $5.4\times10^{-16}$ & $0$ & $2$\\
\midrule
\multirow{2}{*}{強單調上升}
 & 1 & $0$ & $4.6\times10^{-2}$ & $21$\\
 & 2 & $8.7\times10^{-15}$ & $0$ & $22$\\
\bottomrule
\end{tabular}

\vspace{2pt}
\parbox{0.92\textwidth}{\footnotesize 註：$T=24$，權重為 $U(0.2,5)$ 的隨機抽樣，
$h=4$。$z=1$ 時斜率之差與 $\sum_i\mathrm{sign}$ 同時非零，
與命題~\ref{prop:invariance} 的證明一致——該情形下
$t^{\top}(\Delta^{1})^{\top}g=\sum_i g_i$ 不必為零。}
\label{tab:l1inv}
\end{table}

$z=1$ 時斜率不受保護，而其機制可由證明直接讀出：
$\Delta^{1}t=\mathbf 1$，故 $t^{\top}(\Delta^{1})^{\top}g=\sum_i g_i$，
該和不必為零。$z=2$ 時 $\Delta^{2}t=0$，保證因而無條件成立。
\textbf{據此，若要沿時間修勻而保住漂移，$z$ 必須取二階以上， 不可圖省事用一階差分。} 標準 Lee-Carter 的外推之所以沒有這個問題，
是因為它讓漂移自由，零空間因而同為 $\mathrm{span}\{\mathbf 1,t\}$。

\subsection{改變的是解的類別}\label{subsec:l1-sol}

$L_2$ 與 $L_1$ 的差別在解的型態。$L_2$ 的解處處彎曲，
每一個二階差分都非零；$L_1$ 的解為 $z-1$ 次的分段多項式，
折點由資料選出，而 \citet{tibshiraniRJ2014} 證明該估計量在
有界總變差的函數類上達到極小極大最適率。
兩者因而各在不同的函數類上最適：$L_2$ 對應 Sobolev 類（處處平滑），
$L_1$ 對應總變差類（大致平滑而容許少數大的轉折）。

表~\ref{tab:l1kappa} 把兩者調到同一個粗糙度再比較。

\begin{table}[H]
\centering
\caption{臺灣女性 $\kappa_t$ 在同一粗糙度下兩種懲罰的解}
\small
\renewcommand{\arraystretch}{1.1}
\begin{tabular}{llrl}
\toprule
懲罰 & $h$ & 粗糙度 & 非零二階差分（折點）\\
\midrule
原始 $\hat\kappa$ & — & 31.243 & 22 / 22\\
$L_2$（Whittaker） & 0.0420 & 15.604 & 22 / 22\\
$L_1$（trend filtering） & 0.0450 & 15.627 & 15 / 22\\
\bottomrule
\end{tabular}

\vspace{2pt}
\parbox{0.92\textwidth}{\footnotesize 註：2001--2024 年，$z=2$，等權重。
粗糙度為 $\sum_t(\Delta^2\kappa_t)^2$，兩法皆調到原始值的一半。
$L_1$ 的折點位於 2005--2010、2012--2016、2018、2020、2022、2023；
未獲折點者為 2002--2004、2011、2017、2019、2021。
原始 $\hat\kappa$ 二階差分最大的五年為 2015 ($+2.30$)、2010 ($+2.00$)、
2005 ($-1.79$)、2012 ($-1.72$)、2020 ($+1.61$)，五者在 $L_1$ 解中皆保留為折點。}
\label{tab:l1kappa}
\end{table}

在同一粗糙度下，$L_1$ 保留十五個折點而使七段完全筆直，
$L_2$ 則二十二個二階差分全部非零。
稀疏的程度並不大，故不宜誇稱 $L_1$ 給出一條「幾段直線」的路徑；
\textbf{但型態的差別是實質的：$L_1$ 把平滑量集中花在少數幾年， 而原始 $\hat\kappa$ 二階差分最大的五年全部被保留為折點。}
這一點對本文的資料有具體意義——觀測期間含 2021 至 2022 年，
若該期間的衝擊確實與其餘年度不同量級，
則「所有衝擊皆為同一個常態分布」的 $L_2$ 先驗即為誤設，
而 $L_1$ 正是對該誤設的直接回應。

\subsection{貝氏讀法由常態先驗改為尺度混合}\label{subsec:l1-bayes}

$L_1$ 懲罰的貝氏讀法是 Laplace 先驗，而 Laplace 分布是常態的尺度混合，
混合密度為指數分布 \citep{andrews1974}；\citet{parkcasella2008} 即以此
給出 Bayesian lasso 的階層表示。於是

\begin{center}
\begin{minipage}{0.88\textwidth}\small
$L_2$ 懲罰 $\;\longleftrightarrow\;$ 單一變異數的常態先驗 $\;\longleftrightarrow\;$ \textbf{混合模型}\\[3pt]
$L_1$ 懲罰 $\;\longleftrightarrow\;$ 常態的尺度混合先驗 $\;\longleftrightarrow\;$ \textbf{混合分布模型}
\end{minipage}
\end{center}

亦即由 $L_2$ 改為 $L_1$，恰好就是由混合模型走到混合分布模型。
本文先前提出的另一條路線——把穩健損失寫成常態的尺度混合，
使資訊加權成為一個潛在變數層——用的是同一個表示定理
（\citealp{andrews1974}；\citealp{hunter2000}），
只是施加在概似端而非先驗端。兩者因而可以同時進行：
概似端的尺度混合給出穩健與加權，先驗端的尺度混合給出 $L_1$ 修勻。

三項代價須一併記下。其一，$L_1$ 沒有閉式的 $S$，
故後驗變異數 $(W+h\Delta^{\top}\Delta)^{-1}$ 這個現成的區間來源消失；
有效自由度雖可由折點數加 $z$ 計算，但折點是資料選出來的，
以之計算的區間涉及選擇後推論的問題。其二，$h$ 不再是一個變異成分，
故 REML 不適用，須改以交叉驗證或 SURE。其三，沒有 $S$ 也就沒有
命題~\ref{prop:transport} 的輸送算子，故 $L_1$ 不宜用在年齡方向；
但時間方向依推論~\ref{cor:alpha} 本來就不需要輸送，
\textbf{故 $L_1$ 最有吸引力之處恰是它不需要輸送的那個方向。}

最後須明確區隔：此處的 $L_1$ 施加在時間參數的差分上，
目的是選出轉折的位置，與本文先前評估後未採用的
「在 $\beta_x$ 上施以 $L_1$ 以作變數選擇」並非同一件事，
後者的結論不因本節而改變。


\subsection{推廣至收縮先驗}\label{subsec:shrink}

命題~\ref{prop:invariance} 要求懲罰為凸，而該假設其實可以去掉。
文獻在 $L_1$ 之後的發展是把 $\Delta^{z}r$ 的先驗由單一變異數的常態
改為尺度逐點變動的混合：\citet{jullion2007} 以 Normal-Exponential-Gamma、
\citet{faulkner2018} 以 horseshoe 施於 $k$ 階差分並報告其在偏誤與精確度
之間最佳、\citet{kowal2019} 進一步讓局部尺度依賴過程自身的歷史。
$L_1$ 的 Laplace 先驗在這個分類中是尺度混合的一個特例。
以下指出\textbf{整族先驗都保留同樣的兩項不變性}。

\begin{prop}\label{prop:shrink}
設工作概似為 $y\mid r\sim N(r,\,\sigma^2W^{-1})$，
而先驗 $\pi(r)=f(\Delta^{z}r)$ 只透過 $\Delta^{z}r$ 依賴 $r$。
令 $P$ 為對 $\mathcal N=\mathrm{span}\{\mathbf 1,t\}$ 的 $W$-正交投影
（$z\ge2$；$z=1$ 時 $\mathcal N=\mathrm{span}\{\mathbf 1\}$）。則
$$\Ex[Pr\mid y]\;=\;Py ,$$
亦即後驗平均精確保留資料在零空間上的成分。
\end{prop}

\begin{proof}
$\Delta^{z}P=0$，故 $\pi$ 只依賴 $(I-P)r$，
亦即先驗在 $\mathcal N$ 的方向上為平坦。
$P$ 為 $W$-正交投影，故
$(r-y)^{\top}W(r-y)=(Pr-Py)^{\top}W(Pr-Py)
+((I-P)r-(I-P)y)^{\top}W((I-P)r-(I-P)y)$，
交叉項為零。後驗因而在 $Pr$ 與 $(I-P)r$ 之間因式分解，
且 $Pr$ 的後驗為 $N(Py,\sigma^2(P^{\top}WP)^{-1})$，
其平均即 $Py$。
\end{proof}

命題~\ref{prop:shrink} 覆蓋 $L_2$、$L_1$、NEG、horseshoe
與動態收縮，蓋因五者都只透過 $\Delta^{z}r$ 依賴 $r$。
以 horseshoe 的 Gibbs 抽樣核對（$T=24$、保留五萬次、資料含一個轉折）：
加權平均之差為 $7.5\times10^{-5}$、加權斜率之差為 $4.3\times10^{-6}$，
而蒙地卡羅標準誤為 $1.1\times10^{-3}$，故兩者皆落在抽樣誤差之內；
同時後驗確實平滑了資料（$\|\hat r-y\|=1.149$）。

\textbf{其意涵是本文可以採用上述任一先驗而不失去
$\hat\alpha_x$ 與漂移的保證。} 本文選 $L_2$ 與 $L_1$ 的理由因而不是統計效能——
\citet{faulkner2018} 的比較顯示 horseshoe 在偏誤與精確度之間優於
Laplace——而是推導的長度：凸的情形有命題~\ref{prop:invariance} 的
三行次梯度證明，而非凸的情形須走命題~\ref{prop:shrink} 的因式分解論證，
且後者只對高斯工作概似成立。這一點須誠實寫，
不可暗示 $L_1$ 在統計上優於收縮先驗。

\section{高階修正的三種對象}\label{sec:higher}

$b(\mu;c_0)$ 對 Poisson 而言是精確的級數而非展開式，
故「高階」在本文的架構中必須指明對象。以下三種各有不同的理論依據。

\subsection{插入式誤差}\label{subsec:plugin}

本文扣除的是 $b(\hat\mu)$ 而非 $b(\mu)$，其差可展開為

\begin{equation}\label{eq:plugin}
\Ex[b(\hat\mu)]-b(\mu)\;=\;b'(\mu)\,\big(\Ex[\hat\mu]-\mu\big)
\;+\;\tfrac12 b''(\mu)\,\mathrm{Var}(\hat\mu)\;+\;O(\cdot).
\end{equation}

兩個導數皆可由差分恆等式
$\frac{d}{d\mu}\Ex[g(D)]=\Ex[g(D{+}1)]-\Ex[g(D)]$ 精確算出，
施用兩次即得
$b''(\mu)=\Ex[g(D{+}2)]-2\Ex[g(D{+}1)]+\Ex[g(D)]+\mu^{-2}$，
其與數值微分在 $\mu=3$ 處相符至六位（$0.016079$）。
$b''$ 的三項性質：$\mu\to0$ 時 $b''\sim\mu^{-2}$ 而為正且發散
（$\mu=0.02$ 時 $2500.0$）；於 $\mu\approx3.5846$ 變號；
大 $\mu$ 時 $b''\to-\mu^{-3}$，極小值 $-0.0069$ 於 $\mu=4.8575$。

表~\ref{tab:higher} 以 200 次重複檢驗式~\eqref{eq:plugin} 能否解釋實測。

\begin{table}[H]
\centering
\caption{插入式誤差與其二項展開的對照（人數五萬）}
\small
\renewcommand{\arraystretch}{1.1}
\begin{tabular}{lrrrrr}
\toprule
年齡 & $\mu$ & $\Ex[b(\hat\mu)]-b(\mu)$ & 一階項 & 二階項 & 兩項和\\
\midrule
0      & 1.15 & $+0.0020$ & $-0.0247$ & $+0.0211$ & $-0.0036$\\
1--4   & 0.32 & $+0.0006$ & $-0.1132$ & $+0.0754$ & $-0.0378$\\
5--9   & 0.26 & $+0.0075$ & $-0.0862$ & $+0.0829$ & $-0.0032$\\
10--14 & 0.26 & $+0.0493$ & $-0.0298$ & $+0.0746$ & $+0.0447$\\
15--19 & 0.52 & $+0.0447$ & $-0.0071$ & $+0.0362$ & $+0.0290$\\
30--34 & 1.78 & $+0.0016$ & $-0.0100$ & $+0.0098$ & $-0.0002$\\
50--54 & 10.11 & $+0.0009$ & $+0.0014$ & $-0.0005$ & $+0.0009$\\
70--74 & 48.84 & $+0.0004$ & $+0.0008$ & $-0.0002$ & $+0.0005$\\
\bottomrule
\end{tabular}

\vspace{2pt}
\parbox{0.92\textwidth}{\footnotesize 註：重複 200 次，各欄為該年齡各年的平均。
逐年齡相關為 $0.817$、迴歸斜率 $0.792$；一階項與二階項的絕對值
各佔 $47\%$ 與 $53\%$。人數二十萬時相關升至 $0.899$ 而斜率降至 $0.478$，
兩項佔比為 $68\%$ 與 $32\%$。$\hat\mu$ 本身在幼年組偏高 $3\%$ 至 $12\%$。}
\label{tab:higher}
\end{table}

結論有兩層。\textbf{式~\eqref{eq:plugin} 抓到了正確的結構}：逐年齡相關為
$0.817$ 與 $0.899$，故插入式誤差確實由 $\hat\mu$ 的偏誤與變異兩者決定。
\textbf{但它不宜作為校正使用。} 兩項在最關鍵的幼年組幾乎互相抵銷——
$1$--$4$ 歲時一階項為 $-0.1132$、二階項為 $+0.0754$、兩項和 $-0.0378$，
而實際只有 $+0.0006$。亦即兩個各約 $0.1$ 量級的項相減得到 $0.001$ 量級的淨值，
其餘落在三階。要施以二階校正，就得把兩個大項各算到相對精度極高，
而其中的 $\mathrm{Var}(\hat\mu)$ 與 $\Ex[\hat\mu]-\mu$ 本身都只有估計值。
這同時解釋了為何只修勻 $\hat\mu$ 的版本效果不如預期：
修勻降低 $\mathrm{Var}(\hat\mu)$ 而壓低二階項，
卻同時給 $\hat\mu$ 添上平滑偏誤而放大一階項，
而 $|b'|$ 在小 $\mu$ 處甚大（$\mu=0.1$ 時 $-9.31$），後者因而佔上風。

方法論的歸屬清楚。一階的事後修正即 \citet{coxsnell1968}，
其在廣義線性模型的具體形式見 \citet{cordeiro1991}，
而 \citet{kosmidis2009} 推廣至非線性指數族——Lee-Carter 的雙線性結構
正屬此類。更高階的部分，\citet{kosmidis2014} 給出回顧，
而 \citet{kennepagui2017} 區分平均無偏與中位數無偏。
\textbf{後者對本文特別切題，蓋因本文全程以中位數報告結果}，
故若要往高階走，中位數無偏才是前後一致的目標。

\subsection{年齡與時間參數的展開}\label{subsec:beta-order}

$\hat\alpha_x$ 有恆等式而 $\hat\beta_x$ 沒有，原因是前者在奇異值分解之前
即由列平均定出，後者是奇異向量。要寫出 $\Ex[\hat\beta_x]-\beta_x$
須用逐元而非子空間的擾動理論，\citet{abbe2020} 給出的特徵向量逐元展開
正是所需的形式。\textbf{此處的「高階」是指該展開的二階項， 而一階項本身尚未做出}，故這一支的邊際報酬最高，
但也最可能受限於 $A=22$、$T=24$ 的小維度。

\subsection{修勻自身的偏誤}\label{subsec:smooth-order}

第\ref{subsec:nosecond}節的 $(S-I)\log\mu$ 是第三種意思，
而該節的實測與 \citet{calonico2018} 的結論一致：
估計並扣除平滑偏誤不如直接少平滑一些。
三者相較，\ref{subsec:plugin}是解釋性的、\ref{subsec:beta-order}是值得做的、
而\ref{subsec:smooth-order}的理論明白指向不要做。

\section{修勻的先決條件}\label{sec:precond}

前三節把修勻接上中心化校正並量出效果，但未回答一個更前面的問題：
施以修勻的正當性何在，以及該正當性能否事前查核。本節指出兩件事——
條件可以寫成差分算子零空間的一個陳述，而該條件在本文的架構中
恰好由模型自身的假設供給；又該條件所需的量全部有閉式，
因此可以在配適之前算出。

\subsection{偏誤與變異的交換}\label{subsec:tradeoff}

設 $S_h$ 為一族平滑算子且 $S_0=I$，$\hat\theta$ 為待平滑的估計量，
$\Ex[\hat\theta]=\theta+\delta$，$\mathrm{Var}(\hat\theta)=V$。
對任一線性泛函 $c^{\top}$，平滑後的均方誤差為

\begin{equation}\label{eq:mse}
\mathrm{MSE}(h)\;=\;\big(c^{\top}(S_h-I)\theta+c^{\top}S_h\delta\big)^2
\;+\;c^{\top}S_h V S_h^{\top} c .
\end{equation}

\begin{prop}\label{prop:precond}
於 $h=0$ 處對式~\eqref{eq:mse} 微分，得
$$\frac{d}{dh}\mathrm{MSE}(0)\;=\;2\,(c^{\top}\delta)\,
\frac{d}{dh}\big[c^{\top}(S_h\theta+S_h\delta)\big]_{h=0}
\;+\;\frac{d}{dh}\big[c^{\top}S_hVS_h^{\top}c\big]_{h=0}.$$
若 $c^{\top}\delta=0$，則第一項消失，而第二項在平滑降低該泛函變異時為負，
故存在 $h>0$ 使均方誤差嚴格改善。若 $c^{\top}\delta\neq0$，
第一項為一階且其符號取決於資料，改善因而沒有保證。
\end{prop}

\begin{proof}
$S_h\theta-\theta$ 於 $h=0$ 為零，故式~\eqref{eq:mse} 第一項的
平方內部在 $h=0$ 的值為 $c^{\top}\delta$，微分乘積法則即給出上式。
\end{proof}

命題~\ref{prop:precond} 的意涵是次序性的。\textbf{修勻的改善只在它所施加的 估計量對目標泛函已經無偏時才有保證。} 本文的中心化校正所做的正是
把 $\delta$ 移除，於是修勻施加在校正後的估計量上落在有保證的一側；
反之施加在標準 Lee-Carter 上則沒有保證，蓋因該估計量的 $\delta$ 不為零。
第\ref{sec:age}節的實測與此相符：修勻單獨施用於標準 Lee-Carter 時
$\alpha$ 的最大偏誤只由 $0.860$ 動到 $0.714$，
而施加在校正後的估計量上則 $\mathrm{SSE}(\beta)$ 與平均餘命的均方根同時改善。

此處須明確區分兩種說法。命題~\ref{prop:precond} 並未主張修勻施加在
有偏估計量上必定變差，只主張沒有保證；一階項的符號由
$c^{\top}\delta$ 與平滑所致位移的相對方向決定，兩者同向時立刻惡化。

\subsection{零空間即模型}\label{subsec:nullmodel}

$\delta=0$ 之後，剩下的唯一偏誤來源是 $(S_h-I)\theta$，
而該項為零的條件可以寫得很乾淨。

\begin{prop}\label{prop:null}
$z$ 階差分懲罰的零空間為格子指標的 $z-1$ 次以下多項式。
$(S_h-I)\theta=0$ 對一切 $h$ 成立，若且唯若 $\theta$ 落在該零空間內。
\end{prop}

\begin{proof}
$S_h\theta=\theta$ 等價於 $(W+h\Delta^{\top}\!\Delta)\theta=W\theta$，
即 $\Delta^{\top}\!\Delta\theta=0$，而 $W$ 正定且
$\Delta^{\top}\!\Delta$ 半正定，故等價於 $\Delta\theta=0$。
$\Delta^{z}$ 消滅 $z-1$ 次以下的多項式而不消滅其餘，即得。
\end{proof}

命題~\ref{prop:null} 把先決條件由「資料夠平滑嗎」改寫為
「$\theta$ 是否落在懲罰的零空間」，而這是一個關於模型而非關於資料的陳述。
$z=2$ 時零空間為 $\mathrm{span}\{\mathbf 1,x\}$，即\textbf{對數死亡率在
格子上成一次式}。本文的兩個方向各自對應一個已有的死亡率假設。

\begin{itemize}
\item \textbf{年齡方向對應 Gompertz 法則。} $\log\mu_x$ 對年齡成一次式
即 Gompertz 法則的內容，故 $z=2$ 的懲罰在 Gompertz 成立的年齡範圍內
\textbf{沒有平滑偏誤}。
\item \textbf{時間方向對應固定漂移。} $\log\mu_t$ 對時間成一次式
即 $\kappa_t$ 為帶固定漂移的隨機漫步 \citep{lee1992}，
亦即本文外推步驟所用的假設本身。
\end{itemize}

亦即\textbf{先決條件是「懲罰的零空間須包含所相信的模型」}，
而本文兩個方向的結構假設恰好都是 $z=2$ 的零空間。
這同時給出第\ref{sec:time}節那個不變性的另一種讀法：
沿時間修勻之所以不動漂移，與 Lee-Carter 之所以用帶漂移的隨機漫步外推，
是同一件事的兩面。

\subsection{Gompertz 在臺灣資料上成立的範圍}\label{subsec:gompertz}

命題~\ref{prop:null} 把條件化為可查核的量。
表~\ref{tab:curv} 列臺灣女性 2001--2024 年全國配適的
$\Delta^2\log m_x$，以年齡組序為格子。

\begin{table}[H]
\centering
\caption{對數死亡率的年齡曲度}
\small
\renewcommand{\arraystretch}{1.05}
\begin{tabular}{lrr|lrr}
\toprule
年齡 & $\log m$ & $\Delta^2\log m$ & 年齡 & $\log m$ & $\Delta^2\log m$\\
\midrule
0      & $-5.452$ & —        & $50$--$54$ & $-5.932$ & $-0.0023$\\
$1$--$4$   & $-8.351$ & —        & $55$--$59$ & $-5.545$ & $-0.0229$\\
$5$--$9$   & $-9.006$ & $+2.2441$ & $60$--$64$ & $-5.107$ & $+0.0511$\\
$10$--$14$ & $-8.969$ & $+0.6918$ & $65$--$69$ & $-4.627$ & $+0.0416$\\
$15$--$19$ & $-8.302$ & $+0.6304$ & $70$--$74$ & $-4.074$ & $+0.0732$\\
$20$--$24$ & $-8.021$ & $-0.3865$ & $75$--$79$ & $-3.496$ & $+0.0251$\\
$25$--$29$ & $-7.829$ & $-0.0892$ & $80$--$84$ & $-2.880$ & $+0.0386$\\
$30$--$34$ & $-7.521$ & $+0.1166$ & $85$--$89$ & $-2.284$ & $-0.0209$\\
$35$--$39$ & $-7.165$ & $+0.0471$ & $90$--$94$ & $-1.745$ & $-0.0562$\\
$40$--$44$ & $-6.753$ & $+0.0565$ & $95$--$99$ & $-1.298$ & $-0.0931$\\
$45$--$49$ & $-6.341$ & $+0.0002$ & $100^{+}$  & $-0.876$ & $-0.0243$\\
\bottomrule
\end{tabular}

\vspace{2pt}
\parbox{0.92\textwidth}{\footnotesize 註：$5$--$9$ 至 $20$--$24$ 的
$|\Delta^2\log m|$ 平均為 $0.9882$，$25$--$29$ 至 $95$--$99$ 為 $0.1030$，
相差 $9.6$ 倍。$45$--$49$ 至 $55$--$59$ 的曲度為 $0.0002$、$0.0023$、$0.0229$，
亦即 Gompertz 法則在該範圍成立至小數第三位。}
\label{tab:curv}
\end{table}

曲度的分布與人口學的常識一致：$5$--$9$ 至 $20$--$24$ 涵蓋幼年的死亡率谷底
與事故隆起，兩者皆為曲度而非一次式；$25$--$29$ 以上進入 Gompertz 範圍，
曲度降一個數量級。$0$ 與 $1$--$4$ 歲未列曲度，因其組距與其餘不同，
在組序格子上本即不可比——而這兩組正是第\ref{sec:age}節依實作理由
排除在懲罰之外者，此處得到一個不依賴實作的理由。

\subsection{可在配適之前算出的逐年齡診斷}\label{subsec:exante}

把命題~\ref{prop:precond} 與~\ref{prop:null} 合起來，
修勻對 $\hat\alpha_x$ 的均方誤差之比為

\begin{equation}\label{eq:ratio}
\frac{\mathrm{MSE}_{\text{修勻}}(x)}{\mathrm{MSE}_{\text{未修勻}}(x)}
=\frac{\Big(\tfrac1T\sum_t\big[(S_t-I)\log\mu_{\cdot,t}\big]_x\Big)^{2}
+\tfrac1{T^2}\sum_t\big[S_tV_tS_t^{\top}\big]_{xx}}
{\tfrac1{T^2}\sum_t v(\mu_{x,t};c_0)} ,
\end{equation}

其中 $V_t=\mathrm{diag}(v(\mu_{x,t};c_0))$ 為對角——各格的死亡數獨立，
故無須估計共變。\textbf{式~\eqref{eq:ratio} 的每一項都有閉式， 所需的輸入只有曝露數與一組粗略的死亡率，與本文既有的事前可行性檢定相同。}
表~\ref{tab:exante} 以實測中較佳的平滑量計算。

\begin{table}[H]
\centering
\caption{修勻對 $\hat\alpha_x$ 均方誤差的事前判讀}
\small
\renewcommand{\arraystretch}{1.05}
\begin{tabular}{lrrr|lrrr}
\toprule
\multicolumn{4}{c|}{人數五萬，$\lambda=1$} & \multicolumn{4}{c}{人數一萬，$\lambda=3$}\\
年齡 & 平滑偏誤 & $\mathrm{sd}$ 原 & $\sqrt{\text{比}}$ &
年齡 & 平滑偏誤 & $\mathrm{sd}$ 原 & $\sqrt{\text{比}}$\\
\midrule
$5$--$9$   & $-0.0992$ & 0.0707 & 1.638 & $5$--$9$   & $-0.1741$ & 0.0320 & 5.605\\
$10$--$14$ & $+0.1858$ & 0.0716 & 2.704 & $10$--$14$ & $+0.1249$ & 0.0324 & 4.029\\
$15$--$19$ & $-0.1584$ & 0.0978 & 1.705 & $15$--$19$ & $-0.2058$ & 0.0455 & 4.588\\
$20$--$24$ & $-0.1159$ & 0.1189 & 1.069 & $20$--$24$ & $-0.1492$ & 0.0582 & 2.637\\
$25$--$29$ & $+0.0185$ & 0.1324 & \textbf{0.422} & $25$--$29$ & $-0.0006$ & 0.0699 & \textbf{0.565}\\
$40$--$44$ & $+0.0055$ & 0.1102 & \textbf{0.433} & $40$--$44$ & $-0.0005$ & 0.1248 & \textbf{0.435}\\
$55$--$59$ & $+0.0001$ & 0.0562 & \textbf{0.507} & $55$--$59$ & $-0.0004$ & 0.1344 & \textbf{0.417}\\
$70$--$74$ & $+0.0036$ & 0.0303 & \textbf{0.628} & $70$--$74$ & $+0.0156$ & 0.0741 & \textbf{0.580}\\
$85$--$89$ & $-0.0098$ & 0.0241 & \textbf{0.786} & $85$--$89$ & $-0.0240$ & 0.0566 & \textbf{0.731}\\
$100^{+}$  & $+0.0638$ & 0.1086 & \textbf{0.783} & $100^{+}$  & $+0.1159$ & 0.1258 & 1.154\\
\midrule
\multicolumn{3}{l}{全體 $\sqrt{\sum\mathrm{MSE}}$ 之比} & \textbf{0.9607} &
\multicolumn{3}{l}{全體 $\sqrt{\sum\mathrm{MSE}}$ 之比} & 1.0186\\
\bottomrule
\end{tabular}

\vspace{2pt}
\parbox{0.92\textwidth}{\footnotesize 註：比值小於一者修勻改善、大於一者惡化。
為節省篇幅只列代表性年齡，$30$--$39$、$45$--$54$、$60$--$69$、$75$--$84$、
$90$--$99$ 的比值均在 $0.42$ 至 $0.84$ 之間。
$0$ 與 $1$--$4$ 歲在懲罰區外，比值恆為一。}
\label{tab:exante}
\end{table}

三項讀法。其一，\textbf{惡化的年齡恰為表~\ref{tab:curv} 中曲度大的年齡}
（$5$--$9$ 至 $20$--$24$），而 $25$--$29$ 以上均改善且改善幅度可觀——
均方誤差降至四成至八成。亦即式~\eqref{eq:ratio} 的判讀與
命題~\ref{prop:null} 的零空間條件相符，而非只是一個經驗規則。

其二，\textbf{該診斷能預測實測中依目標而異的結果}。人數五萬時全體比值為
$0.9607$（改善），人數一萬時為 $1.0186$（略為惡化），
而第\ref{sec:age}節的實測正是：人數一萬時 $\lambda$ 加大使
$\mathrm{SSE}(\beta)$ 與平均餘命改善，但 $\alpha$ 的中位偏誤由
$0.0471$ 升到 $0.0963$。診斷的對象是 $\hat\alpha_x$，故它預測的是後者。
\textbf{平滑量不存在對三個參數同時最適的單一值}，這一點在選 $\lambda$ 時
須明寫。

其三，$100^{+}$ 在人數一萬時惡化。該組是格子的端點，差分懲罰在端點的
自由度最少，而晚年死亡率的減速又使其偏離 Gompertz。
第\ref{sec:next}節所列的組距加權差分算子不處理這一項，
端點的處理須另行考慮。

\subsection{零格與補值：三個層級}\label{subsec:tiers}

$5$--$9$ 至 $20$--$24$ 與 $0$、$1$--$4$ 不滿足零空間條件，
而這些年齡恰好也是零格最集中之處，故替代作法須與零格的處理一併考慮。

一個零格所能借到的資訊有三個來源：秩一結構（經 $\beta_x\kappa_t$ 連到
同年齡的其他年度與同年度的其他年齡）、鄰近年齡（修勻），
以及參考母體（標準化死亡比）。零格替代值 $c_0$ 是其中的退化情形，
它不向任何地方借，只填一個常數。三者的適用範圍由兩個量決定，
而兩個量都算得出來。

第一個是\textbf{列資訊量} $I_x=\sum_t\mu_{x,t}$。Poisson 分布對
$\log\mu$ 的 Fisher 資訊為 $\mu$，故 $I_x$ 即該年齡全期對
$\alpha_x$ 的資訊總量。它同時給出整列全零的機率 $e^{-I_x}$
與期望零格數 $\sum_t e^{-\mu_{x,t}}$。
第二個是\textbf{局部曲度} $|\Delta^2\log\mu_x|$，即命題~\ref{prop:null}
的零空間條件。據此分三級。

\begin{table}[H]
\centering
\caption{零格的分布與三個層級（臺灣女性年齡結構）}
\small
\renewcommand{\arraystretch}{1.05}
\begin{tabular}{lrrrrl}
\toprule
年齡 & $I_x$ & 期望零格數 & $\Pr(\text{整列全零})$ & $|\Delta^2\log m|$ & 層級\\
\midrule
0          & 5.54 & 19.1 & $3.9\times10^{-3}$ & — & 丙\\
$1$--$4$   & 1.51 & 22.5 & $2.2\times10^{-1}$ & — & 丙\\
$5$--$9$   & 1.23 & 22.8 & $2.9\times10^{-1}$ & 2.2441 & 丙\\
$10$--$14$ & 1.26 & 22.8 & $2.8\times10^{-1}$ & 0.6918 & 丙\\
$15$--$19$ & 2.50 & 21.6 & $8.2\times10^{-2}$ & 0.6304 & 丙\\
$20$--$24$ & 4.11 & 20.2 & $1.6\times10^{-2}$ & 0.3865 & 丙\\
$25$--$29$ & 5.99 & 18.7 & $2.5\times10^{-3}$ & 0.0892 & 乙\\
$30$--$34$ & 8.54 & 16.8 & $2.0\times10^{-4}$ & 0.1166 & 乙\\
$35$--$39$ & 12.61 & 14.2 & $3.3\times10^{-6}$ & 0.0471 & 乙\\
$40$--$44$ & 23.79 & 8.9 & $4.6\times10^{-11}$ & 0.0565 & 甲\\
$45$--$49$ & 35.12 & 5.6 & $5.6\times10^{-16}$ & 0.0002 & 甲\\
$55$--$59$ & 72.59 & 1.3 & $3.0\times10^{-32}$ & 0.0229 & 甲\\
$70$--$74$ & 234.41 & 0.0 & $1.6\times10^{-102}$ & 0.0732 & 甲\\
$85$--$89$ & 354.79 & 0.0 & $8.3\times10^{-155}$ & 0.0209 & 甲\\
$100^{+}$  & 24.27 & 8.7 & $2.9\times10^{-11}$ & 0.0243 & 甲\\
\bottomrule
\end{tabular}

\vspace{2pt}
\parbox{0.92\textwidth}{\footnotesize 註：人數一萬，觀測 24 年故每列 24 格。
為節省篇幅省略 $50$--$54$、$60$--$69$、$75$--$84$、$90$--$99$，
其 $I_x$ 均在 $48$ 以上而期望零格數在 $3.2$ 以下。
分級規則為 $I_x\ge20$ 列甲，否則 $|\Delta^2\log m|<0.3$ 列乙、其餘列丙。
期望零格總數為 $207.1$ 格，佔全部 $528$ 格的 $39.2\%$；
人數五千與兩千時為 $48.7\%$ 與 $63.0\%$。}
\label{tab:tiers}
\end{table}

\begin{itemize}
\item \textbf{甲級（$I_x$ 充足）。} 零格雖有但列資訊量足夠，
中心化校正即足，不需借用外部資訊。
\item \textbf{乙級（$I_x$ 不足而曲度小）。} 零空間條件成立，
故可沿年齡修勻借用鄰近年齡——此即第\ref{subsec:exante}節判為改善的範圍。
\item \textbf{丙級（$I_x$ 不足且曲度大）。} 修勻在此既無資訊可借
又會引入偏誤，須改用參考母體，或改以不依賴多項式零空間的約束。
\end{itemize}

丙級的替代工具有三。其一是參考母體，即 \citet{lee2003} 的 partial SMR
與 \citet{yue2019} 的接法；其代價是本文已指出的——型態不符時
$\alpha$ 偏誤惡化至九至十五倍。其二是改以單調性或凸性的約束取代平滑懲罰，
此即 \citet{broffitt1988} 在貝氏修勻一側所處理的問題；
該類約束的可行集不是多項式空間，故不受命題~\ref{prop:null} 的限制，
而嬰幼兒死亡率雖陡降卻確實單調，約束因而成立。
\textbf{這一項不需要參考母體，故不受資料取得的牽制，是三者中最值得先試的。}
其三是沿時間而非沿年齡借資訊——依第\ref{sec:time}節，
時間方向的零空間條件由本文自己的漂移假設供給，
故在這些年齡仍然成立；但整列全零時無效（期望有 $0.90$ 個年齡如此）。
其四是跨地區借資訊，即把固定年齡、各鄉鎮市區的死亡數視為一個
Poisson 複合決策問題，以非參數經驗貝氏估出死亡率的分布
\citep{favaro2024}。該法在零格上不需任何替代值，蓋因 $y=0$ 時
Robbins 形式 $p_G(1)/p_G(0)$ 仍為正。初步評估見
\texttt{研究筆記/複合決策與經驗貝氏\_評估.md}：
丙級年齡的對數死亡率均方誤差可由 $5.54$ 降至 $0.115$（oracle 為 $0.063$），
但須把混合分布放在率而非計數均值上，而該變體超出所引文獻的理論範圍。

丙級恰為 $0$ 至 $20$--$24$ 的六組，\textbf{而這與第\ref{subsec:exante}節
由均方誤差比判為惡化的年齡完全相同}——兩者所用的量不同
（一為列資訊量與曲度，一為平滑偏誤與解析變異），結論一致，
故該分界不是任一條規則的產物。這同時解釋了實測中殘餘的 $\alpha$ 最大偏誤
為何固定落在 $1$--$4$ 歲且不隨平滑量變動：該組是丙級，
而丙級的病灶在定義上就不是修勻所能處理的。

整列全零的期望年齡數 $\sum_x e^{-I_x}$ 亦為閉式，
其值在人數兩千、五千、一萬時為 $4.26$、$2.09$、$0.90$，
而 \texttt{規模極限\_待討論.md} 以模擬量到的對應值為
$4.3$、$1.9$、$1.0$。\textbf{該欄原以模擬取得，此處指出它有閉式}，
故三個層級的判定全部可在配適之前完成。

\subsection{修勻與補值的三種結合，其理論依據並不相同}\label{subsec:combine}

\textbf{其一，補值在前、修勻在後。} 先把零格填為
$\mathrm{SMR}_t\cdot m^{R}_{x,t}$ 再修勻與校正。
此法有一項容易忽略的要求：\textbf{補值後該格已不是 Poisson 觀測，
而是參考母體資料的確定性函數}，故 $b(\mu;c_0)$ 不是它的偏誤、
$v(\mu;c_0)$ 不是它的變異。
$b$ 是 $\log\max(D,c_0)$ 的偏誤，而補值後的量是 $\mathrm{SMR}\cdot m^R$，
兩者的分布無關。於是\textbf{校正只能施於未補值的格，
而補值格的權重須另由參考母體的抽樣變異給出}。
本文現行的「零格以 SMR 填補」只做配適而未接校正，故尚無此問題；
但下一步若要把兩者合併，這一項必須先處理。

\textbf{其二，修勻與補值同時，即懲罰化 Poisson 概似。}

\begin{equation}\label{eq:penpois}
\max_{\theta}\ \sum_{x,t}\big[D_{x,t}\,\theta_{x,t}-E_{x,t}e^{\theta_{x,t}}\big]
\;-\;h\,\|\Delta^{z}\theta\|^2 .
\end{equation}

式~\eqref{eq:penpois} 中零格的貢獻為 $-E e^{\theta}$，
完全不需要取零的對數，\textbf{故 $c_0$ 與補值都不出現——
沒有東西需要補}。這是三者之中理論上最乾淨的結合，
而它已經存在：\citet{delwarde2007} 把差分懲罰直接加在 Poisson 對數概似上，
\citet{currie2013} 給出帶約束的一般框架，\citet{brouhns2002} 則是
Poisson 對數雙線性本身。

但它與本文的主軸有一處衝突，須明白寫出。
式~\eqref{eq:penpois} 迴避了對數轉換的偏誤而非校正它，
依 Firth 自己的用語屬\textbf{預先}而非\textbf{事後}的修正，
與 Firth 同屬一格；而本文已誠實記錄 Firth 在 $\alpha$ 的最大偏誤上全面勝出。
故式~\eqref{eq:penpois} 的定位應為對照組的上界參考，
與 partial SMR 在本文中的地位相同，而非本文的提案。

\textbf{其三，以自我一致性把兩者接成同一個不動點。}
本文的校正值 $\ell_{x,t}-b(\hat\mu_{x,t})$ 依建構其期望恰為
$\log\mu_{x,t}$，\textbf{故它本身即是一個自我一致的補值，
只是施於全部的格而非只施於零格}。
把現行的迭代（配適 $\to$ 算 $\hat\mu$ $\to$ 扣 $b$ $\to$ 重配適）
看成一個不動點，修勻便只是改動其中的配適一步，
三者因而收在同一條估計方程上。
須誠實記下的是，扣 $b$ 並不是任何概似的 E 步，
故這只是與懲罰化 EM \citep{dempster1977} 結構平行，而非其實例；
單調收斂須另行證明。實測的收斂在三至六次迭代之間，
但那是實證結果。

\textbf{三者相較}：其二理論最乾淨但定位是對照組；
其一可行且與第零點五優先的計畫一致，但須先處理補值格的分布問題；
其三是唯一能把本文既有的校正與修勻寫成單一估計方程的路線，
故理論上的新意最多，代價是收斂性尚待證明。

\section{混合模型讀法的前提}\label{sec:premise}

式~\eqref{eq:unified} 的推論讀法要成立，以下各項須先確認。
前兩項是代數上的，其餘為資料上的。

\begin{enumerate}[label=\arabic*.]
\item \textbf{差分的階數決定漂移是否受罰。} $\Delta^2$ 的零空間為
$\mathrm{span}\{\mathbf 1,t\}$，故 $z\ge2$ 不罰水準與斜率；
純 $z=1$ 會縮減斜率（核對中由 $0.2508$ 縮到 $0.2433$）。
\item \textbf{漂移估計式須與懲罰相容。} 現行的端點差不在零空間內，
應改為加權最小平方斜率，見第\ref{subsec:time-caveat}節。
\item \textbf{變異成分須可識別。} 實測的適當 $\lambda$ 隨人數而變
（一萬約 $0.3$、五萬約 $0.1$），故固定一個值不可行；
但 REML 可能把變異成分推到邊界，而該情形下的概似比檢定是非標準的
\citep{crainiceanu2004}。
\item \textbf{常態近似在時間方向較安全。} $\hat\kappa_t$ 是跨二十二個年齡組的
匯總，有中央極限定理可用；而 $\hat\alpha_x$ 是單一年齡沿時間的平均，
$\mu\ll1$ 時 $\log\max(D,c_0)$ 是有原子的離散分布。
可疑的是年齡方向而非時間方向。
\item \textbf{權重應計算而非估計。} 混合模型要求誤差變異已知至一個尺度，
而本文已有 $v(\mu;c_0)=\mathrm{Var}[\log\max(D,c_0)]$ 的閉式，
直接代入即可。
\item \textbf{衝擊的序列獨立性須檢查。} 高斯差分先驗假設 $\Delta^{z}\kappa$
為白噪音。觀測期間含 2021 至 2022 年，若該期間的衝擊有序列相關，
先驗即為誤設，而平滑量會以錯誤的方式吸收它。
\item \textbf{整列全零仍無法處理。} 修勻算子精確保留常數（$S\mathbf 1=\mathbf 1$），
故某年齡各年全為零死亡時，沿時間修勻無從增加資訊。
此與年齡方向在 $1$--$4$ 歲的盲點同性質。
\end{enumerate}

\section{後續工作}\label{sec:next}

依本文的結果，優先序如下。

\textbf{第一，把差分算子改為組距加權。} 現行作法把 $0$ 與 $1$--$4$ 歲排除在
懲罰之外，而排除本身製造了一個人為邊界：人數五萬、$\lambda=10$ 時
$\alpha$ 最大偏誤移到 $5$--$9$ 歲且符號相反（$-0.303$，
$\lambda=1$ 時為 $10$--$14$ 歲的 $+0.277$）。改以組距加權的差分算子
可使二十二組落在同一個懲罰內，既不誤設不等距的格子，也不製造邊界。

\textbf{第二，把命題~\ref{prop:transport} 與命題~\ref{prop:invariance} 接進主報告。} 兩者都是短證明且已有實測支持，
而命題~\ref{prop:invariance} 的一般形式使 $L_2$ 與 $L_1$ 共用同一組保證。

\textbf{第三，取得 \citet{kimeldorf1967} 與 \citet{whittaker1923} 的原文。}
兩篇皆已在報告中被引用並據以實作，卻尚未持有全文；
前者是第\ref{subsec:one}節貝氏讀法的原始依據。

\textbf{第四，$1$--$4$ 歲改以參考母體或沿時間借資訊。}
第\ref{sec:premise}節第七項指出修勻在該年齡無從施力，
而 \citet{lee2003} 的 partial SMR 與 \citet{yue2019} 的接法正是替代工具。

\textbf{第五，尺度混合的階層模型與其後驗區間。}
概似端與先驗端的尺度混合可同時進行，而後驗給出本文目前所缺的區間推論；
計算須走 Laplace 近似 \citep{rue2009} 而非抽樣。

\clearpage
\bibliographystyle{plainnat}
\bibliography{../../報告/references,references_theory1}

\end{document}
